En cada apartado se propone una representación de una función o de una curva y una EDO. Decide si representa una solución real en todo el dominio indicado. Si hay que restringir el dominio o separar ramas, indica los intervalos abiertos máximos en los que se obtiene una solución. Justifica cada decisión mediante derivación y comprueba la regularidad necesaria; una igualdad que sólo se cumple en puntos aislados no basta.
\begin{enumerate}
\item $y=\ln(1+x^2)$, $x\in\mathbb{R}$; \quad $(1+x^2)y'=2x$.
\item $y=\dfrac{x}{1+x^2}$, $x\in\mathbb{R}$; \quad $(1+x^2)y'=1$.
\item $y=|x|$, $x\in\mathbb{R}$; \quad $(y')^2=1$.
\item $y^3+y=x+\sin x$, con $x,y\in\mathbb{R}$; \quad $(1+3y^2)y'=1+\cos x$. Justifica también si la relación determina una única función de $x$.
\item $y^2=\ln(1+x)$, con $x>-1$; \quad $2(1+x)yy'=1$.
\item $x=t+t^3$, $y=t^2$, $t\in\mathbb{R}$; \quad $(1+3y)^2(y')^2=4y$.
\item $x=t^3$, $y=t^6+t^3$, $t\in\mathbb{R}$; \quad $y'=2x+1$. Examina expresamente qué ocurre en $t=0$.
\end{enumerate}
